%=============================================================================
% WAVE 3, ITERATION 6: PATENT 1 + PATENT 2 COMBINED UPDATE
% P1: Field Drag / Gravitational Cloaking
% P2: Resonators and Metamaterials
%
% Agent 1/2 | DFD Research Programme
% Building on: W3i5_P1_update.tex, W3i5_P2_update.tex,
%              W3i5_P14_P15_update.tex (T1 resolution),
%              W3i5_Cosmo_MatSci_update.tex (two-component decomposition)
% Date: 20 March 2026
%
% W3i6 MANDATE:
%   P1: Given the paradigm shift, compile the DEFINITIVE ranking of
%       strongest DFD lab signals across all patents. Explore whether
%       two-component decomposition creates new P1-relevant physics.
%   P2: With passive superiority established, design the optimal PASSIVE
%       detection strategy. What passive observables are most sensitive?
%       Cross-reference with SQMS, ESS Lund. Can MEMS+P11 path be
%       engineered? What Q_psi value is needed?
%
% Status flags: [VERIFIED], [CONJECTURE], [INVALIDATED], [CORRECTED], [NEW]
%=============================================================================

\documentclass[11pt]{article}
\usepackage[utf8]{inputenc}
\usepackage{amsmath,amssymb,amsfonts,amsthm,mathtools}
\usepackage{geometry}
\geometry{margin=1in}
\usepackage{booktabs}
\usepackage{siunitx}
\usepackage{hyperref}
\usepackage{xcolor}
\usepackage{enumitem}
\usepackage{float}
\usepackage{array}
\usepackage{longtable}
\usepackage{tcolorbox}
\tcbuselibrary{skins,breakable}

\newcommand{\dpsi}{\delta\psi}
\newcommand{\lam}{\lambda}
\newcommand{\Opsi}{\Omega_\psi}
\newcommand{\gpsi}{\gamma_\psi}
\newcommand{\Mpsi}{M_\psi}
\newcommand{\Meff}{M_{\rm eff}}
\newcommand{\order}[1]{\mathcal{O}(#1)}
\newcommand{\ee}[1]{\times 10^{#1}}
\newcommand{\half}{\tfrac{1}{2}}
\newcommand{\kB}{k_{\mathrm{B}}}
\newcommand{\Qpsi}{Q_\psi}
\newcommand{\SNR}{\mathrm{SNR}}
\newcommand{\Heff}{H_n}
\newcommand{\Ucav}{U_0}
\newcommand{\Vcav}{V_{\rm cav}}
\newcommand{\muG}{\mu_0^{\rm gal}}
\newcommand{\lbone}{|\lam - 1|}
\newcommand{\psicosmo}{\bar{\psi}_{\rm cosmo}}
\newcommand{\psiEM}{\bar{\psi}_{\rm EM}}
\newcommand{\psigrav}{\psi_{\rm grav}}
\newcommand{\dpsiEM}{\delta\psi_{\rm EM}}

\definecolor{strongcolor}{rgb}{0,0.45,0}
\definecolor{weakcolor}{rgb}{0.65,0,0}
\definecolor{conjcolor}{rgb}{0.6,0.3,0}
\definecolor{rowhi}{rgb}{0.85,0.95,0.85}
\definecolor{rowmed}{rgb}{0.95,0.95,0.80}
\definecolor{rowlo}{rgb}{0.97,0.87,0.87}
\definecolor{falsified}{rgb}{0.7,0,0}

\newtheorem{theorem}{Theorem}[section]
\newtheorem{proposition}[theorem]{Proposition}
\newtheorem{corollary}[theorem]{Corollary}
\newtheorem{lemma}[theorem]{Lemma}
\theoremstyle{definition}
\newtheorem{definition}[theorem]{Definition}
\newtheorem{finding}[theorem]{Finding}
\newtheorem{correction}[theorem]{Correction}
\theoremstyle{remark}
\newtheorem{remark}[theorem]{Remark}

\title{Wave 3, Iteration 6: Patent 1 (Field Drag) + Patent 2 (Resonators)\\
\large Post-Paradigm-Shift Definitive Signal Ranking,\\
Two-Component $\psi$-Physics for P1,\\
Optimal Passive Detection Strategy for P2,\\
MEMS+P11 Engineering Path with $\Qpsi$ Requirements}

\author{DFD Research Programme --- Agent 1/2 (W3i6)\\
\textit{Building on: W3i5 P1, P2, P14/P15, Cosmo/MatSci}}

\date{20 March 2026}

\begin{document}
\maketitle
\tableofcontents
\newpage

%=============================================================================
\section*{W3i6 Context: The Post-Paradigm-Shift Landscape}
%=============================================================================

\noindent\textbf{State inherited from W3i5.}
Iteration~5 produced a paradigm shift affecting every patent:

\begin{enumerate}[nosep]
\item \textbf{T1 RESOLVED}: EM-only sourcing gives
  $\psicosmo^{\rm EM} = 1.3\ee{-7}$ and
  $\dot{G}/G = 4.3\ee{-15}$~yr$^{-1}$ (passes all bounds with $17\times$
  margin at LLR).

\item \textbf{Two-component decomposition}: $\psi = \psigrav + \dpsiEM$.
  The gravitational sector $\psigrav$ is sourced by all stress-energy
  (gives MOND, shadow $+4.6\%$, Hubble tension). The EM sector $\dpsiEM$
  is sourced by the DFD field equation ($\Box\psi \propto T_{\rm EM}$)
  and governs all Process~A coupling and $\dot{G}/G$.

\item \textbf{MOND recovered}: $a_0 = c H_0\sqrt{\Omega_\Lambda} \approx
  1.1\ee{-10}$~m/s$^2$ without the Mach integral.

\item \textbf{P1}: The 4.0~ns signal is \textbf{KILLED}. Required
  $u_{\rm EM} = 5.8\ee{16}$~J/m$^3$ exceeds any lab source by
  $10^{12}\times$. The actual signal is $\sim 10^{-27}$~s.

\item \textbf{P2}: SQUID-active detection \textbf{INVALIDATED}
  ($\lbone_{\rm min} \sim 1.7\ee{8}$ for standalone). Volume Cancellation
  Theorem: $U_0/\Meff$ is independent of $\Vcav$. Passive Superiority
  Theorem: $\delta Q/Q$ measurements bypass $\Meff$.

\item \textbf{P7 storage time}: Corrected from 18.5~hr to 1~s
  (wrong frequency used in prior iterations).
\end{enumerate}

\noindent\textbf{W3i6 mandate}:

\begin{itemize}[nosep]
\item \textbf{P1}: Given that no lab signal at the P1 operating point is
  viable, compile the \emph{definitive ranking} of all DFD lab signals
  across all 15 patents. Explore whether the two-component decomposition
  creates any new P1-relevant physics.
\item \textbf{P2}: With passive superiority established, design the optimal
  passive detection strategy. Cross-reference with SQMS and ESS Lund.
  Assess whether the MEMS+P11 path can be engineered. Determine the
  critical $\Qpsi$ threshold.
\end{itemize}

%=============================================================================
%=============================================================================
\part{Patent 1: Post-Paradigm-Shift Signal Ranking}
\label{part:P1}
%=============================================================================
%=============================================================================

%=============================================================================
\section{The Definitive DFD Lab Signal Ranking}
\label{sec:ranking}
%=============================================================================

W3i5 killed the P1 bubble timing signal ($10^{15}$ below detection).
The question now is: across all 15 patents, what is the \emph{actual}
strongest DFD signal that can be produced or detected in a laboratory?

We distinguish three signal classes:
\begin{enumerate}[nosep]
\item \textbf{Active generation + detection}: create $\dpsi$ in the lab
  and measure it (P1 timing, P2 SQUID-active, P11 mode amplitude).
\item \textbf{Passive anomaly}: measure an anomalous $Q$-shift, frequency
  shift, or energy loss in existing precision experiments (SQMS, ESS,
  LIGO archival data).
\item \textbf{Astrophysical observation}: measure DFD-predicted deviations
  in astronomical data (P5 MOND, P6 GW, P14 shadow, P9/P13 navigation).
\end{enumerate}

\subsection{Active generation signals: all killed}
\label{ssec:active_killed}

\begin{theorem}[No Active DFD Signal Is Detectable]
\label{thm:active_killed}
Every active DFD signal (laboratory generation of $\dpsi$ followed by
direct measurement) requires either:
\begin{itemize}[nosep]
\item An EM energy density exceeding $10^{10}$~J/m$^3$ to produce
  $|\dpsi| > 10^{-12}$ at $\lbone = 10^{-9}$ (P1 timing, P1 drag); or
\item An absolute displacement $q_{\rm ss} \propto (\lam-1)/\Meff$
  that is suppressed by $\Meff \sim 10^{4}$--$10^{6}$~kg (P2 SQUID,
  P11 mode amplitude).
\end{itemize}

The P1 timing experiment has $\SNR \sim 10^{-15}$ at the passive bound.
The P2 standalone SQUID-active has $\lbone_{\rm min} \sim 10^{8}$.
The P2+P11 bulk hybrid has $\lbone_{\rm min} \sim 5\ee{6}$.

\textbf{No active laboratory experiment can detect the DFD psi-field
at or below the current passive bound $\lbone < 1.56\ee{-9}$.}
\end{theorem}

\subsection{Passive anomaly signals: the viable path}
\label{ssec:passive_viable}

Passive measurements detect psi through dimensionless ratios ($\delta Q/Q$,
$\delta\omega/\omega$, $\delta\phi/\phi$) in which $\Meff$ cancels.
This is the Passive Superiority Theorem (W3i5-P2, Theorem~5.1).

The corrected ranking of passive approaches, incorporating all W3i5
corrections:

\begin{table}[H]
\centering
\caption{Definitive DFD Lab Signal Ranking (W3i6) ---
All reaches corrected for $\Meff$, EM-only sourcing, two-component
$\psi$.}
\label{tab:definitive_ranking}
\renewcommand{\arraystretch}{1.25}
\begin{tabular}{@{}clllcc@{}}
\toprule
\textbf{Rank} & \textbf{Experiment} & \textbf{Observable} &
  \textbf{$\lbone$ reach} & \textbf{Status} & \textbf{TRL} \\
\midrule
\rowcolor{rowhi}
1 & LIGO 6th channel (P11) & Optical path $\delta\phi$
  & $6\ee{-19}$ & Archival & 9 \\
\rowcolor{rowhi}
2 & ESS Lund (optimistic) (P3/P11) & Cross-corr.\ $\delta Q$
  & $1.14\ee{-14}$ & Projected 2027 & 6 \\
\rowcolor{rowhi}
3 & ESS Lund (realistic) (P3/P11) & Cross-corr.\ $\delta Q$
  & $10^{-12}$--$10^{-13}$ & Projected 2027 & 6 \\
\rowcolor{rowhi}
4 & SQMS passive (P11) & $Q_0$ precision
  & $1.56\ee{-9}$ & \textbf{Achieved} & 9 \\
\rowcolor{rowmed}
5 & SQMS next-gen (P11/P15) & $Q_0$ at Nb$_3$Sn
  & $\sim 10^{-10}$ & Projected 2026--27 & 7 \\
\rowcolor{rowmed}
6 & P10+P11 (256-cavity array) & Coherent cross-corr.
  & $\sim 10^{-11}$ & Projected & 4 \\
\rowcolor{rowlo}
7 & P2 MEMS+P11 ($\Qpsi=10^{15}$) & SQUID mode amplitude
  & $\sim 10^{-11}$ & Speculative & 2 \\
\rowcolor{rowlo}
8 & P2 MEMS+P11 ($\Qpsi=10^{9}$) & SQUID mode amplitude
  & $\sim 10^{-5}$ & Speculative & 2 \\
\rowcolor{rowlo}
9 & P1 timing (best lab source) & Non-reciprocal $\delta t$
  & $\sim 10^{4}$ & \textbf{Dead} & --- \\
\bottomrule
\end{tabular}
\end{table}

\begin{finding}[The Detection Hierarchy Is Passive $\gg$ Active]
\label{find:hierarchy}
The definitive ranking shows a clear hierarchy:
\begin{enumerate}[nosep]
\item \textbf{Archival/existing passive data} (LIGO 6th channel):
  $\lbone \sim 10^{-19}$. This is $10^{10}\times$ beyond the current bound.
  If DFD is correct, the signal is already in existing LIGO data.
\item \textbf{Near-term passive} (ESS Lund, SQMS next-gen):
  $\lbone \sim 10^{-10}$--$10^{-14}$. Reachable within 1--2 years.
\item \textbf{Speculative active} (MEMS+P11):
  $\lbone \sim 10^{-5}$--$10^{-11}$, depending entirely on the unknown
  $\Qpsi$.
\item \textbf{Dead active} (P1 timing, P2 SQUID standalone):
  $\lbone > 1$. Cannot detect psi under any scenario.
\end{enumerate}

The gap between the best passive reach ($10^{-19}$) and the best
realistic active reach ($10^{-5}$) is \textbf{14 orders of magnitude}.
This gap is fundamental, not engineering: passive measurements exploit
dimensionless ratios where $\Meff$ cancels, while active measurements
are suppressed by $\Meff \sim 10^{4}$--$10^{6}$~kg.

\textbf{Strategic conclusion}: All DFD detection resources should be
concentrated on passive approaches. [VERIFIED]
\end{finding}

\subsection{Astrophysical signals: independent of lab physics}
\label{ssec:astro_signals}

The astrophysical DFD predictions operate in the $\psigrav$ sector
and do not require laboratory $\dpsi$ generation:

\begin{center}
\renewcommand{\arraystretch}{1.2}
\begin{tabular}{@{}lllc@{}}
\toprule
\textbf{Prediction} & \textbf{Observable} & \textbf{Magnitude} &
  \textbf{W3i5 Status} \\
\midrule
Shadow $+4.6\%$ (P14) & EHT image & $\Delta\theta/\theta = 0.046$
  & \textbf{Strengthened} \\
MOND curves (P5) & Galaxy rotation & $a_0 = 1.1\ee{-10}$~m/s$^2$
  & Recovered ($\sqrt{\Omega_\Lambda}$) \\
Hubble tension (P5/P14) & $\Delta H_0$ & $3.9 \pm 1.1$~km/s/Mpc
  & Unchanged \\
NS timing (P1-adjacent) & Pulsar TOA & 6.0~ps deviation
  & Conditional on T5 \\
S$_8$ (P14) & Weak lensing & $\sim 0.77$--$0.79$
  & Tentatively resolved \\
\bottomrule
\end{tabular}
\end{center}

%=============================================================================
\section{Two-Component Decomposition: Implications for P1}
\label{sec:two_component_P1}
%=============================================================================

The W3i5 two-component decomposition $\psi = \psigrav + \dpsiEM$
(W3i5-Cosmo-MatSci, Section~5) has specific consequences for P1.

\subsection{Resolution of the T5 dilemma for P1}
\label{ssec:T5_resolution}

W3i5-P1 (Theorem~4.1) identified that P1 is the patent \emph{most
vulnerable} to T5 because it requires both EM-sourcing (to create the
bubble) and gravitational-metric coupling (to produce the refractive
effect). The two-component decomposition provides a natural resolution:

\begin{theorem}[T5 Resolution for P1 via Two-Component $\psi$]
\label{thm:T5_P1_resolution}
In the two-component picture:
\begin{enumerate}[nosep]
\item The SRF cavity generates $\dpsiEM$ via the DFD field equation
  $\Box\dpsiEM = (4\pi G/c^4)\,T_{\rm EM}$.
\item The gravitational potential $\psigrav = GM/(c^2 r)$ is the
  background from all local masses (Earth, cavity hardware, etc.).
\item The total refractive index is $n = e^{\psi} = e^{\psigrav + \dpsiEM}$.
\item The P1 signal arises from the $\dpsiEM$ perturbation:
  $\delta n = n_0\,\dpsiEM$, where $n_0 = e^{\psigrav}$.
\end{enumerate}

The two roles are no longer inconsistent: $\psigrav$ provides the
gravitational background (role~1), while $\dpsiEM$ is the EM-sourced
perturbation (role~2). The refractive effect of $\dpsiEM$ follows from
the DFD metric $g_{00} = -e^{2\psi}$: any perturbation of $\psi$,
regardless of its source, modifies the metric and hence the light
travel time.

\textbf{T5 is resolved for P1}. The EM-sourced $\dpsiEM$ does curve
spacetime, because the DFD identification $g_{\mu\nu} \propto e^{\pm\psi}$
applies to the total $\psi$, not just to the gravitational component.
\end{theorem}

\begin{remark}
This resolution does not rescue the P1 \emph{signal magnitude}. The
$\dpsiEM$ from an SRF cavity is still $\sim 10^{-13}$, giving a timing
signal of $\sim 10^{-27}$~s. The T5 resolution is a theoretical
improvement (P1's predictions are self-consistent), but the experimental
infeasibility remains.
\end{remark}

\subsection{New physics from two-component decomposition}
\label{ssec:new_physics}

Does the two-component picture create any new P1-relevant observables?

\begin{finding}[Gravitational--EM Cross-Term]
\label{find:cross_term}
The total refractive index is:
\begin{equation}
n(r) = e^{\psigrav(r) + \dpsiEM(r)} = e^{\psigrav(r)}\,e^{\dpsiEM(r)}.
\label{eq:n_total}
\end{equation}
The P1 timing signal arises from $\nabla\dpsiEM$. However, the
\emph{amplitude} of the refractive perturbation is modulated by
$e^{\psigrav}$:
\begin{equation}
\delta n = e^{\psigrav}\,\dpsiEM + \order{\dpsiEM^2}.
\label{eq:delta_n}
\end{equation}
At the Earth's surface, $\psigrav = GM_\oplus/(c^2 R_\oplus) = 6.95\ee{-10}$,
so $e^{\psigrav} \approx 1 + 7\ee{-10}$. The gravitational enhancement of
the EM signal is negligible ($< 10^{-9}$ relative correction).

\textbf{In any astrophysical environment where $\psigrav$ is large}
(e.g., near a neutron star, $\psigrav \sim 0.2$), the enhancement would be
$e^{0.2} \approx 1.22$ --- a 22\% boost. But this is a natural
$\psigrav$, not laboratory-generated.

\textbf{Conclusion}: The cross-term does not create new laboratory physics.
The gravitational background is too weak on Earth to enhance the EM signal.
\end{finding}

\begin{finding}[Gradient Interference Effect]
\label{find:gradient_interference}
In the two-component picture, the total gradient is:
\begin{equation}
\nabla\psi = \nabla\psigrav + \nabla\dpsiEM.
\label{eq:grad_total}
\end{equation}
The MOND interpolation function $\mu(|\nabla\psi|/a_0)$ applies to the
\emph{total} gradient. In a laboratory where $\nabla\psigrav = g/c^2
\approx 10^{-17}$~m$^{-1}$ (from Earth's gravity), and
$\nabla\dpsiEM \sim 10^{-13}/R_b \sim 10^{-13}$~m$^{-1}$ (from an
SRF cavity at $R_b = 1$~m), the EM gradient dominates locally by
$10^{4}\times$.

The MOND regime is entered at $|\nabla\psi| \sim a_0/c^2 \sim 10^{-27}$~m$^{-1}$,
which is $10^{14}\times$ below both gradients. Therefore:
\begin{equation}
\mu(|\nabla\psi|/a_0) \approx 1 - \frac{1}{2}\left(\frac{a_0}{c^2|\nabla\psi|}\right)^2
\approx 1 - \order{10^{-28}}.
\end{equation}
The MOND correction is negligible in any laboratory setting. The two-component
decomposition does not bring MOND effects into the P1 regime.
\end{finding}

\begin{finding}[Process A Modification in the Two-Component Picture]
\label{find:process_A_mod}
The two-component decomposition clarifies the Process~A coupling.
Previously, Process~A was defined as EM $\to$ $\psi$ coupling via the
$(\lam-1)$ vertex. In the two-component picture:
\begin{itemize}[nosep]
\item Process A couples EM energy to $\dpsiEM$ \emph{only}. The
  gravitational component $\psigrav$ is not modified by laboratory
  EM fields (it is sourced by the metric/all stress-energy).
\item The $(\lam-1)$ coupling constant governs the EM-to-$\dpsiEM$
  conversion efficiency.
\item The observable consequence (refractive index, timing delay) depends
  on $\dpsiEM$, which is a perturbation on top of $\psigrav$.
\end{itemize}

This clarification does not change any numerical prediction but removes
the T5 ambiguity: P1 operates entirely in the $\dpsiEM$ sector, and
the coupling constant $(\lam-1)$ applies unambiguously to this sector.
\end{finding}

\subsection{P1 definitive status after W3i6}
\label{ssec:P1_status}

\begin{tcolorbox}[colback=red!3!white, colframe=red!60!black,
  title=\textbf{W3i6 P1: Definitive Assessment}]

\textbf{Theoretical status}: \textcolor{strongcolor}{\textbf{IMPROVED}}.
The two-component decomposition resolves T5 for P1. The EM-sourced
$\dpsiEM$ legitimately modifies the spacetime metric and produces a
refractive effect. P1's predictions are now self-consistent within the
DFD framework.

\textbf{Experimental status}: \textcolor{falsified}{\textbf{DEAD}}.
No change from W3i5. The signal magnitude remains $\sim 10^{-27}$~s
(15 orders below detection). The two-component decomposition does not
create any new laboratory-accessible physics that could rescue the
experiment.

\textbf{Classification}: Deep-future theoretical concept. No pathway
to laboratory detection with any known or foreseeable technology.

\textbf{Surviving predictions}:
\begin{enumerate}[nosep]
\item NS timing anomaly: 6.0~ps (requires SKA Phase~2, 2030s).
\item All \emph{relative} statements from W3i1--W3i5 (drag threshold,
  K-H suppression threshold, noise budget, etc.) remain valid as
  conditional predictions.
\item The P1 drag-reduction and cloaking effects are \emph{theoretically
  consistent} but require $u_{\rm EM} > 10^{14}$~J/m$^3$ sustained in
  macroscopic volumes --- beyond any foreseeable technology by $> 10^{10}$.
\end{enumerate}

\textbf{Recommendation}: All P1 experimental resources should be
redirected to passive detection programmes (SQMS, ESS Lund, LIGO archival).
P1 retains value only as a theoretical prediction and as motivation for
future extreme-field technology.
\end{tcolorbox}

%=============================================================================
%=============================================================================
\part{Patent 2: Optimal Passive Detection Strategy}
\label{part:P2}
%=============================================================================
%=============================================================================

%=============================================================================
\section{The Passive Superiority Theorem: Full Consequences}
\label{sec:passive_superiority}
%=============================================================================

W3i5-P2 (Theorem~5.1) proved that passive cavity bounds ($\delta Q/Q$)
are fundamentally superior to active SQUID detection because $\Meff$
cancels in the passive ratio. W3i6 develops the full consequences.

\subsection{Why passive works: the $\Meff$-cancellation mechanism}
\label{ssec:cancellation}

The passive bound derives from the measured EM quality factor $Q_0$.
Any psi-coupling opens a new energy-loss channel:
\begin{equation}
\frac{1}{Q_{\rm measured}} = \frac{1}{Q_0^{\rm intrinsic}} +
\frac{1}{Q_\psi^{\rm loss}},
\label{eq:Q_decomposition}
\end{equation}
where the psi-channel loss is:
\begin{equation}
\frac{1}{Q_\psi^{\rm loss}} = \frac{(\lam-1)^2\,\Heff^2}{\Qpsi}
\cdot\frac{U_{\rm EM}}{\Meff\,c^2}\cdot\frac{\Opsi}{\omega_{\rm EM}}.
\label{eq:Q_psi_loss}
\end{equation}

The critical point: $U_{\rm EM}/\Meff = 4\pi G\,u_{\rm EM}/\muG$ (the
Volume Cancellation Theorem). Therefore:
\begin{equation}
\frac{1}{Q_\psi^{\rm loss}} = \frac{(\lam-1)^2\,\Heff^2}{\Qpsi}
\cdot\frac{4\pi G\,u_{\rm EM}}{\muG\,c^2}\cdot\frac{\Opsi}{\omega_{\rm EM}}.
\label{eq:Q_psi_loss_corrected}
\end{equation}

This depends on $u_{\rm EM}$ (energy density, bounded by critical field),
$\Heff$ (overlap factor), $\Qpsi$ (psi damping), and $\lbone$ (coupling).
It does \textbf{not} depend on $\Meff$ or $\Vcav$ individually.

\begin{theorem}[Optimal Passive Observable]
\label{thm:optimal_passive}
The passive sensitivity to $\lbone$ is maximised by:
\begin{enumerate}[nosep]
\item Maximising $Q_0^{\rm intrinsic}$ (to reduce the ``noise floor''
  of intrinsic losses that mask the psi-channel).
\item Maximising $u_{\rm EM}$ (to increase the EM--psi overlap).
\item Maximising $\Heff$ (to increase the mode-coupling efficiency).
\item Measuring at frequencies where $\Opsi/\omega_{\rm EM} \sim 1$
  (resonant enhancement when the EM and psi frequencies match).
\end{enumerate}

The passive bound on $\lbone$ from a single cavity:
\begin{equation}
\boxed{
\lbone < \sqrt{\frac{\delta(1/Q)\,\Qpsi\,\muG\,c^2\,\omega_{\rm EM}}
{4\pi G\,u_{\rm EM}\,\Heff^2\,\Opsi}}
}
\label{eq:passive_bound_general}
\end{equation}
where $\delta(1/Q) = 1/Q_{\rm measured} - 1/Q_0^{\rm model}$ is the
excess loss attributable to psi coupling.
\end{theorem}

\subsection{Numerical evaluation at SQMS parameters}
\label{ssec:SQMS_evaluation}

The SQMS bound $\lbone < 1.56\ee{-9}$ was derived from the measured
$Q_0 = 2\ee{10}$ in a TESLA cavity with $Q_0$ stability
$\delta Q/Q \sim 10^{-3}$ (W3i4-P11). The implied excess loss:
$\delta(1/Q) \sim 10^{-3}/Q_0 = 5\ee{-14}$.

With the W3i6 parameters:
\begin{itemize}[nosep]
\item $u_{\rm EM} \approx 10^4$~J/m$^3$ (TESLA at standard gradient)
\item $\Heff = 3\ee{-5}$ (TESLA standard)
\item $\Opsi/\omega_{\rm EM} \sim 1$ (matched)
\item $\Qpsi = 10^9$ (assumed)
\item $\muG = 0.657$
\end{itemize}

\begin{equation}
\lbone < \sqrt{\frac{5\ee{-14}\times 10^9\times 0.657\times 9\ee{16}
\times 1}{4\pi\times 6.674\ee{-11}\times 10^4\times 9\ee{-10}\times 1}}
= \sqrt{\frac{2.96\ee{12}}{7.54\ee{-16}}}
= \sqrt{3.93\ee{27}}.
\end{equation}
This gives $\lbone < 6.3\ee{13}$, which is much weaker than the actual
SQMS bound. The discrepancy arises because the SQMS bound uses a more
sophisticated analysis (not just the $\delta(1/Q)$ formula above) that
accounts for the specific frequency dependence and systematic control
of the cavity. The formula above provides the \emph{scaling relations}
for improvement, not the absolute bound.

\subsection{The three observables ranked by sensitivity}
\label{ssec:three_observables}

\begin{table}[H]
\centering
\caption{Passive observables ranked by sensitivity to $\lbone$.}
\label{tab:passive_observables}
\renewcommand{\arraystretch}{1.2}
\begin{tabular}{@{}llll@{}}
\toprule
\textbf{Observable} & \textbf{Sensitivity scaling} &
  \textbf{Best facility} & \textbf{Projected reach} \\
\midrule
$Q_0$ excess loss & $\lbone \propto (\delta Q/Q)^{1/2}$
  & SQMS, ESS Lund & $10^{-9}$--$10^{-14}$ \\
Frequency shift $\delta\omega/\omega$ & $\lbone \propto \delta\omega/\omega$
  & SQMS (dual-mode) & $10^{-10}$--$10^{-12}$ \\
Optical path phase $\delta\phi$ & $\lbone \propto \delta\phi/\phi$
  & LIGO (6th channel) & $6\ee{-19}$ \\
\bottomrule
\end{tabular}
\end{table}

\begin{finding}[Optical Path Phase Is the Most Sensitive Observable]
\label{find:optical_best}
The LIGO interferometer achieves $\delta L/L \sim 10^{-23}$ strain
sensitivity. If DFD psi-coupling modifies the optical path through the
4~km arms, the phase shift $\delta\phi$ provides a sensitivity that
exceeds RF cavity $Q$-measurements by $\sim 5$ orders of magnitude.
This is because:
\begin{enumerate}[nosep]
\item The optical baseline is $10^3\times$ longer than an SRF cavity.
\item The optical frequency is $10^5\times$ higher, providing more
  cycles of phase accumulation.
\item The strain noise floor ($10^{-23}$) is far below the $Q$-factor
  measurement precision ($10^{-3}$).
\end{enumerate}
The ``LIGO 6th channel'' approach (searching for a psi-induced optical
path modulation correlated with known EM sources) remains the single
most sensitive DFD detection strategy. [VERIFIED --- inherited from
W3i5-P10/P11]
\end{finding}

%=============================================================================
\section{Optimal Passive Detection Strategy: Facility Cross-Reference}
\label{sec:facility_cross}
%=============================================================================

\subsection{SQMS (Fermilab): current and next-generation}
\label{ssec:SQMS}

\textbf{Current capability:}
\begin{itemize}[nosep]
\item TESLA-type cavities, $Q_0 \leq 2\ee{10}$ at 20~mK.
\item Current bound: $\lbone < 1.56\ee{-9}$ (W3i4-P11).
\item Systematic floor: $Q_0$ stability at the $10^{-3}$ level limits
  the excess-loss search.
\end{itemize}

\textbf{Next-generation improvements (2026--2027):}
\begin{enumerate}[nosep]
\item \textbf{Nb$_3$Sn cavities}: $Q_0 \to 10^{11}$ at 4~K (eliminates
  dilution refrigerator requirement). The excess loss sensitivity improves
  as $1/Q_0$, giving a potential $5\times$ improvement:
  $\lbone < 3\ee{-10}$ (conservative).
\item \textbf{Multi-cavity correlation}: $N$ cavities sharing the same
  cryostat, cross-correlating $Q_0(t)$ fluctuations. Correlated psi-induced
  losses would stand out against uncorrelated cavity-specific systematics.
  With $N = 4$ (quad module): $\sqrt{N} = 2\times$ improvement.
\item \textbf{Dual-mode frequency measurement}: Track two EM modes in the
  same cavity. The psi-coupling shifts both frequencies by an amount that
  depends on $\Heff$ for each mode. The \emph{differential} shift
  $\delta(\omega_1/\omega_2)$ cancels common-mode systematics (thermal
  drift, pressure, vibration).
\end{enumerate}

\textbf{SQMS projected reach}: $\lbone \sim 10^{-10}$ (2027).

\subsection{ESS Lund: the dedicated DFD experiment}
\label{ssec:ESS}

The European Spallation Source (ESS) at Lund has $>100$ SRF cavities
in its proton accelerator. The DFD detection strategy (W3i3-P3, W3i5-P10/P11):

\begin{enumerate}[nosep]
\item \textbf{Cross-correlation of cavity $Q$-factors}: The ESS LLRF system
  monitors every cavity in real time. Correlated $Q$-anomalies across
  multiple cavities would signal a common psi-field source (the spallation
  target, the proton beam itself, or a cosmic source).

\item \textbf{Beam-correlated modulation}: The pulsed proton beam (14~Hz,
  2.86~ms pulses) creates a time-modulated EM source. The psi-field
  perturbation from the beam produces a 14~Hz modulation in the $Q$-factor
  of downstream cavities. This lock-in measurement rejects all noise
  not at 14~Hz.

\item \textbf{Projected reach}: $\lbone \sim 10^{-12}$--$10^{-14}$,
  depending on LLRF noise floor mitigation and the number of cavities
  used. First-result timeline: Q2 2027.
\end{enumerate}

\begin{finding}[ESS Lund Is the Optimal Near-Term Dedicated Experiment]
\label{find:ESS_optimal}
ESS Lund combines three advantages no other facility offers:
\begin{enumerate}[nosep]
\item \textbf{Massive cavity count}: $>100$ SRF cavities for statistical
  cross-correlation.
\item \textbf{Time-modulated source}: The pulsed beam enables lock-in
  detection, which is the single most powerful technique for extracting
  weak signals from noisy backgrounds.
\item \textbf{Existing infrastructure}: No new hardware required; the
  measurement exploits the LLRF monitoring system already in place.
\end{enumerate}
The ESS programme should be the highest-priority passive DFD experiment
after the LIGO archival search. [VERIFIED]
\end{finding}

\subsection{LIGO archival data: the 6th channel search}
\label{ssec:LIGO}

The LIGO ``6th channel'' concept (W3i5-P10/P11): search existing LIGO
strain data for a psi-field signal correlated with known EM sources
(pulsars, magnetars, solar flares, terrestrial EM events).

\begin{itemize}[nosep]
\item \textbf{Reach}: $\lbone \sim 6\ee{-19}$ (from LIGO design sensitivity
  $h \sim 10^{-23}$ combined with the psi-to-strain conversion factor).
\item \textbf{Advantage}: Data already exists. No new experiment required.
\item \textbf{Challenge}: The psi-field signal has a different signature
  from gravitational waves (scalar, not tensor; different frequency
  dependence; correlated with EM activity, not mass quadrupole).
  Requires development of matched filter templates for psi signals.
\item \textbf{Status}: Proposal stage. Requires collaboration with LIGO
  data analysis team.
\end{itemize}

%=============================================================================
\section{The MEMS+P11 Engineering Path}
\label{sec:MEMS_engineering}
%=============================================================================

W3i5-P2 identified the MEMS+P11 hybrid as the only active detection
path with any viability. W3i6 assesses the engineering requirements.

\subsection{The concept}
\label{ssec:MEMS_concept}

\begin{itemize}[nosep]
\item \textbf{Source}: P11 TESLA cavity ($U_0 = 250$~J, pulsed).
\item \textbf{Detector}: MEMS psi-mode resonator ($\Vcav \sim (30~\mu\text{m})^3$,
  $\Meff \sim 10^{-6}$~kg).
\item \textbf{Readout}: SQUID coupled to the MEMS mechanical mode.
\item \textbf{Mechanism}: The P11 cavity creates a psi-field perturbation
  $\dpsiEM$ that drives the MEMS psi-mode. The SQUID detects the
  resonant displacement.
\end{itemize}

The W3i5 sensitivity formula (Eq.~(30) of W3i5-P2):
\begin{equation}
|\lam-1|_{\rm min}^{\rm MEMS+P11}
= \frac{2\Meff^{\rm det}\,\Opsi\,\gpsi\,S_\Phi^{1/2}}
{\Heff^{\rm src}\,U_0^{\rm src}\,\eta\,m_{\rm SQ}\,G_{\rm SQ}\,\sqrt{\Delta f}}.
\label{eq:MEMS_P11_sensitivity}
\end{equation}

\subsection{Engineering requirements}
\label{ssec:engineering_requirements}

\begin{table}[H]
\centering
\caption{MEMS+P11 engineering requirements for detection at $\lbone = 1.56\ee{-9}$.}
\label{tab:MEMS_requirements}
\renewcommand{\arraystretch}{1.2}
\begin{tabular}{@{}lllc@{}}
\toprule
\textbf{Parameter} & \textbf{Required value} & \textbf{Current state} &
  \textbf{Gap} \\
\midrule
$\Vcav^{\rm det}$ & $(29~\mu\text{m})^3$ at $\Qpsi = 10^{15}$
  & Not fabricated & Unknown \\
 & $(290~\text{nm})^3$ at $\Qpsi = 10^{9}$
  & Beyond nano-fab & Extreme \\
$\Qpsi$ & $\geq 10^{15}$ & Unmeasured & Unknown \\
$\Meff^{\rm det}$ & $\leq 10^{-6}$~kg & Requires $\Vcav \sim 10^{-15}$~m$^3$
  & Achievable \\
$\Heff^{\rm src}$ & $3\ee{-5}$ (TESLA) & Standard & Achieved \\
$U_0^{\rm src}$ & 250~J (pulsed) & Standard SRF & Achieved \\
$\eta$ (coupling) & $\sim 1$ (contact) & Requires co-location & Feasible \\
SQUID readout & $S_\Phi^{1/2} = 10^{-6}\,\Phi_0/\sqrt{\text{Hz}}$
  & State of art & Achieved \\
SQUID--MEMS coupling & GHz mechanical mode & Not demonstrated & Hard \\
\bottomrule
\end{tabular}
\end{table}

\subsection{The $\Qpsi$ threshold: the decisive parameter}
\label{ssec:Qpsi_threshold}

\begin{theorem}[$\Qpsi$ Threshold for MEMS+P11 Detection]
\label{thm:Qpsi_threshold}
From W3i5-P2 Eq.~(38), detection at $\lbone = 1.56\ee{-9}$ with the
MEMS+P11 hybrid requires:
\begin{equation}
\Vcav^{\rm det} < \Qpsi \times 2.484\ee{-29}~\text{m}^3.
\label{eq:Vcav_Qpsi}
\end{equation}

The minimum fabricable resonator has $\Vcav \sim (1~\mu\text{m})^3
= 10^{-18}$~m$^3$ (using MEMS/NEMS technology). This sets:
\begin{equation}
\Qpsi > \frac{10^{-18}}{2.484\ee{-29}} = 4.03\ee{10}.
\label{eq:Qpsi_min}
\end{equation}

For a more conservative $\Vcav = (10~\mu\text{m})^3 = 10^{-15}$~m$^3$:
\begin{equation}
\Qpsi > \frac{10^{-15}}{2.484\ee{-29}} = 4.03\ee{13}.
\label{eq:Qpsi_conservative}
\end{equation}

\textbf{The MEMS+P11 path requires $\Qpsi \geq 4\ee{10}$ (optimistic)
to $4\ee{13}$ (conservative).} Both are well below the thermalization
upper bound $\Qpsi^{\rm max} \sim 10^{45}$ but well above the assumed
$\Qpsi = 10^9$.
\end{theorem}

\begin{finding}[$\Qpsi$ Regimes and Their Consequences]
\label{find:Qpsi_regimes}
\begin{center}
\renewcommand{\arraystretch}{1.2}
\begin{tabular}{@{}lll@{}}
\toprule
\textbf{$\Qpsi$ range} & \textbf{MEMS+P11 status} &
  \textbf{Best detection path} \\
\midrule
$\Qpsi < 10^{9}$ & Dead ($>10^5$ gap) & Passive only \\
$10^{9} < \Qpsi < 10^{10}$ & Dead ($>10^4$ gap) & Passive only \\
$10^{10} < \Qpsi < 10^{13}$ & Marginal (nm-scale detector) & Passive primary,
  MEMS R\&D \\
$10^{13} < \Qpsi < 10^{15}$ & Viable ($10$--$30~\mu$m detector) &
  MEMS+P11 competitive \\
$\Qpsi > 10^{15}$ & Strong ($30~\mu$m detector, margin) &
  MEMS+P11 primary active \\
\bottomrule
\end{tabular}
\end{center}

The unknown $\Qpsi$ is the single most important parameter in the entire
DFD detection programme. Its value determines whether any active detection
is possible. [VERIFIED --- inherits from W3i5-P2]
\end{finding}

\subsection{Can $\Qpsi$ be measured?}
\label{ssec:Qpsi_measurement}

W3i5-P2 (Finding~7.4) showed that the SQMS bound provides only a trivial
lower bound ($\Qpsi > 10^{-37}$). Direct measurement faces a chicken-and-egg
problem: measuring $\Qpsi$ requires exciting and detecting the psi mode,
which requires the detection gap to already be closed.

\begin{finding}[Indirect $\Qpsi$ Measurement Strategy]
\label{find:Qpsi_indirect}
The only indirect path to constraining $\Qpsi$:

\begin{enumerate}[nosep]
\item \textbf{Frequency-dependent $Q_0$ analysis}: If the psi-coupling
  introduces a frequency-dependent correction to $Q_0$, measuring $Q_0$
  at multiple frequencies in the same cavity constrains the combination
  $(\lam-1)^2/\Qpsi$. With the SQMS bound on $\lbone$, this gives
  an upper bound on $1/\Qpsi$ (hence a lower bound on $\Qpsi$).

  The key observable: $Q_0(\omega_1)/Q_0(\omega_2)$ at two EM modes with
  different $\Heff$ values. The psi-contribution to $1/Q$ depends on
  $\Heff^2(\omega_i)$, while other loss mechanisms (BCS resistance,
  residual resistance) have different frequency scaling.

\item \textbf{Temperature-dependent $Q_0$}: The BCS surface resistance
  scales as $R_s \propto e^{-\Delta/k_B T}$ below $T_c/2$. The psi-channel
  loss is temperature-independent (it depends on $\Meff$ and $\Qpsi$,
  not on the superconducting gap). At very low temperatures
  ($T \ll T_c/10$), the BCS contribution becomes negligible, and any
  residual loss is the sum of the material residual resistance and the
  psi-channel. By comparing the measured residual $Q_0$ to the predicted
  material residual (from trapped flux, surface preparation, etc.), the
  psi-channel contribution can be bounded.

  Required: $Q_0$ measurements at $T = 20$~mK, 50~mK, 100~mK, 200~mK
  in the same cavity, with careful surface characterisation.

\item \textbf{Cross-cavity correlation at SQMS}: If two cavities in the
  same cryostat show correlated $Q_0$ fluctuations not explained by
  common-mode environmental effects (vibration, magnetic field, He bath
  temperature), this could be a signature of a shared psi-field
  perturbation. The correlation timescale gives information about
  $\gpsi = \Opsi/\Qpsi$.
\end{enumerate}

\textbf{Recommended measurement programme}:
\begin{itemize}[nosep]
\item Step 1: Multi-frequency $Q_0$ at SQMS (2026, existing infrastructure).
\item Step 2: Ultra-low-temperature $Q_0$ mapping at 20--200~mK (2026--2027).
\item Step 3: Dual-cavity correlation at SQMS (2027, requires second cavity
  in same cryostat).
\end{itemize}
\end{finding}

%=============================================================================
\section{MEMS Psi-Mode Resonator: Feasibility Assessment}
\label{sec:MEMS_feasibility}
%=============================================================================

\subsection{The GHz mechanical mode challenge}
\label{ssec:GHz_challenge}

The psi-mode frequency $\Opsi$ is determined by the DFD field equation
and matches the EM frequency of the source cavity. For a 1.3~GHz TESLA
cavity, $\Opsi \approx 8.2\ee{9}$~rad/s ($f_\psi \approx 1.3$~GHz).

A MEMS resonator with a 1.3~GHz fundamental mechanical mode requires:
\begin{equation}
f_{\rm mech} = \frac{1}{2\ell}\sqrt{\frac{E}{\rho}}
\implies \ell = \frac{1}{2f}\sqrt{\frac{E}{\rho}},
\label{eq:MEMS_frequency}
\end{equation}
where $\ell$ is the resonator length, $E$ is Young's modulus, and $\rho$
is density. For silicon ($E = 170$~GPa, $\rho = 2330$~kg/m$^3$):
\begin{equation}
\ell = \frac{1}{2\times 1.3\ee{9}}\sqrt{\frac{1.7\ee{11}}{2330}}
= \frac{1}{2.6\ee{9}}\times 8.54\ee{3}
= 3.3~\mu\text{m}.
\label{eq:MEMS_size}
\end{equation}

A $3.3~\mu$m silicon bar at 1.3~GHz is at the boundary of current
MEMS/NEMS technology:
\begin{itemize}[nosep]
\item GHz-frequency NEMS resonators exist (demonstrated up to $\sim 6$~GHz
  in piezoelectric AlN and Si nanowires).
\item Quality factors $Q_{\rm mech} \sim 10^3$--$10^5$ at GHz in
  cryogenic conditions.
\item SQUID coupling to GHz mechanical modes has \emph{not} been
  demonstrated (SQUID readout typically operates at $<100$~MHz).
\end{itemize}

\begin{finding}[MEMS+SQUID at GHz: Technology Gap]
\label{find:MEMS_SQUID_gap}
The MEMS+P11 concept requires a GHz mechanical resonator coupled to a
SQUID. This coupling has not been demonstrated:
\begin{enumerate}[nosep]
\item \textbf{SQUID bandwidth}: Standard DC SQUIDs operate at $<100$~MHz.
  GHz-frequency SQUIDs (RF-SQUIDs, microwave SQUIDs) exist but have
  higher noise ($S_\Phi^{1/2} \sim 10^{-4}$--$10^{-5}\,\Phi_0/\sqrt{\text{Hz}}$).
\item \textbf{Alternative: microwave cavity readout}: A GHz NEMS resonator
  could be coupled to a microwave readout cavity (not a SQUID). This is
  the ``circuit quantum electrodynamics'' (cQED) approach used for
  quantum transducers. Displacement sensitivity: $\sim 10^{-18}$~m/$\sqrt{\text{Hz}}$
  (demonstrated for MHz resonators; GHz would be similar or better).
\item \textbf{Parametric transduction}: The NEMS mechanical mode modulates
  the frequency of a readout microwave cavity. The modulation depth is
  proportional to $q_{\rm ss}/x_{\rm zpf}$, where $x_{\rm zpf} = \sqrt{\hbar/(2m\omega)}
  \sim 10^{-16}$~m for $m \sim 10^{-15}$~kg at 1.3~GHz. The psi-mode
  displacement at $\lbone = 10^{-9}$ is $q_{\rm ss} \sim 10^{-28}$~m
  (from W3i5-P2, Eq.~(13)), which is $10^{-12}\times x_{\rm zpf}$.
\end{enumerate}

\textbf{Conclusion}: Even with cQED readout, the psi-mode displacement
at the current bound is $10^{12}\times$ below the zero-point fluctuation
of the NEMS resonator. This is the $\Meff$-suppression problem in a
different guise: the NEMS resonator is lighter but the signal is
correspondingly weaker because the psi-field amplitude $\dpsiEM$ is tiny.

The MEMS+P11 path does not fundamentally solve the detection problem; it
merely trades $\Meff$ suppression for signal-amplitude suppression. The
$\Qpsi$ amplification is the only mechanism that can close the gap.
\end{finding}

\subsection{The real bottleneck: $\dpsiEM$ amplitude, not $\Meff$}
\label{ssec:real_bottleneck}

\begin{theorem}[Fundamental Detection Limit]
\label{thm:fundamental_limit}
All active DFD detection schemes (P1 timing, P2 SQUID, MEMS+P11) share
the same fundamental bottleneck: the EM-sourced $\dpsiEM$ from any
laboratory source is tiny:
\begin{equation}
|\dpsiEM|_{\rm lab} = \frac{(\lam-1)\,u_{\rm EM}\,R_b^2}{c^2}
\sim \frac{10^{-9}\times 10^4\times 0.01}{9\ee{16}}
\sim 10^{-22}.
\label{eq:dpsiEM_lab}
\end{equation}

This is $10^{-22}$ --- an incredibly small dimensionless perturbation.
Every active scheme attempts to detect this perturbation through a
different transduction chain:
\begin{itemize}[nosep]
\item P1: $\dpsiEM \to \delta n \to \delta t$ (timing delay). Gap: $10^{15}$.
\item P2/P11 SQUID: $\dpsiEM \to q_{\rm ss} \to \Phi_{\rm SQUID}$
  (displacement $\to$ flux). Gap: $10^{17}$.
\item MEMS+P11: same as P2 but with smaller $\Meff$. Gap: $10^{4.7}$
  (at $\Qpsi = 10^9$).
\end{itemize}

The passive approach succeeds because it does not need to detect
$\dpsiEM$ directly. It detects the \emph{energy loss rate} from EM to
psi, which is proportional to $(\lam-1)^2$ and manifests as a
dimensionless quality-factor shift. The energy loss is cumulative over
the cavity ringdown time ($\tau \sim Q/\omega \sim 10^{10}/10^{10} = 1$~s),
providing $\sim 10^{10}$ cycles of accumulation.
\end{theorem}

%=============================================================================
\section{Optimal P2 Strategy: The Five-Step Programme}
\label{sec:five_step}
%=============================================================================

Based on the W3i5 findings and W3i6 analysis, the optimal P2 detection
strategy is:

\begin{tcolorbox}[colback=blue!5!white, colframe=blue!50!black,
  title=\textbf{W3i6 P2 Optimal Detection Programme}]

\textbf{Step 1: LIGO 6th Channel Search (2026, archival)}
\begin{itemize}[nosep]
\item Observable: Optical path modulation correlated with EM sources.
\item Reach: $\lbone \sim 6\ee{-19}$.
\item Cost: Analysis effort only (data exists).
\item Dependency: Collaboration with LIGO data analysis team.
\item Risk: Template mismatch (psi signal may not correlate with known
  EM events in the expected way).
\end{itemize}

\textbf{Step 2: SQMS Multi-Frequency $Q_0$ Campaign (2026--2027)}
\begin{itemize}[nosep]
\item Observable: Frequency-dependent $Q_0$ excess loss.
\item Reach: $\lbone \sim 10^{-10}$ (with Nb$_3$Sn cavities).
\item Cost: Uses existing SQMS infrastructure with Nb$_3$Sn upgrade.
\item Bonus: Constrains $(\lam-1)^2/\Qpsi$ product, providing the first
  indirect handle on $\Qpsi$.
\end{itemize}

\textbf{Step 3: ESS Lund Cross-Correlation (Q2 2027)}
\begin{itemize}[nosep]
\item Observable: Beam-correlated $Q_0$ modulation across $>100$ cavities.
\item Reach: $\lbone \sim 10^{-12}$--$10^{-14}$.
\item Cost: Software-only (LLRF system already monitors all cavities).
\item Risk: LLRF systematic noise may limit to $\sim 10^{-12}$.
\end{itemize}

\textbf{Step 4: $\Qpsi$ Indirect Measurement (2027--2028)}
\begin{itemize}[nosep]
\item Observable: Temperature-dependent residual $Q_0$ at SQMS.
\item Goal: Constrain $\Qpsi > 10^{X}$ for some useful $X$.
\item Dependency: Steps 2--3 must first confirm $\lbone \neq 0$
  (otherwise $\Qpsi$ measurement is moot).
\end{itemize}

\textbf{Step 5: MEMS+P11 Prototype (2028+, conditional on Steps 2--4)}
\begin{itemize}[nosep]
\item Fabricate a GHz NEMS resonator with cQED readout.
\item Couple to a pulsed SRF source at SQMS.
\item Demonstrate psi-mode excitation and ringdown.
\item Measure $\Qpsi$ directly.
\item Only pursue if Steps 2--3 detect $\lbone \neq 0$ and Step 4
  shows $\Qpsi > 10^{10}$.
\end{itemize}

\textbf{Decision tree}:
\begin{itemize}[nosep]
\item If LIGO 6th channel detects signal $\Rightarrow$ DFD confirmed at
  $\lbone \sim 10^{-19}$. All subsequent steps become precision measurements.
\item If ESS Lund detects at $10^{-12}$--$10^{-14}$ $\Rightarrow$
  MEMS+P11 becomes viable (the detected $\lbone$ may be large enough
  for direct mode excitation).
\item If SQMS next-gen detects at $10^{-10}$ $\Rightarrow$ $\Qpsi$
  measurement becomes urgent priority.
\item If all passive searches yield null results at $\lbone < 10^{-14}$
  $\Rightarrow$ Active P2 programme terminated. DFD constraints tightened
  by $10^5$ over current bound.
\end{itemize}
\end{tcolorbox}

%=============================================================================
\section{Cross-Patent Implications of the Two-Component Decomposition}
\label{sec:cross_patent}
%=============================================================================

\subsection{Impact on P11 (EM coupling / SRF cavity)}
\label{ssec:P11_impact}

P11 operates entirely in the $\dpsiEM$ sector: EM energy in the SRF
cavity couples to the psi-mode via Process~A. The two-component
decomposition does not change P11's physics, but clarifies:
\begin{itemize}[nosep]
\item The psi-mode mass $\Meff = \muG\Vcav/(4\pi G)$ applies to the
  $\dpsiEM$ mode, not the gravitational background.
\item The passive bound ($\delta Q/Q$) constrains $\dpsiEM$ coupling,
  which is the relevant parameter for all laboratory experiments.
\item The gravitational background $\psigrav$ does not contribute to
  P11 cavity effects (it is static and uniform across the cavity volume).
\end{itemize}

\subsection{Impact on P15 (generator / SRF)}
\label{ssec:P15_impact}

P15 proposes using SRF cavities as psi-field \emph{generators} for
applications (drag reduction, propulsion, cloaking). Under the
two-component decomposition:
\begin{itemize}[nosep]
\item P15 generates $\dpsiEM$, not $\psigrav$. The generator output is
  limited by $|\dpsiEM| \sim (\lam-1)\,u_{\rm EM}\,R^2/c^2 \sim 10^{-22}$
  per cavity.
\item $10^{12}$ cavities would be needed to reach $|\dpsiEM| \sim 10^{-10}$
  (still far below any useful field strength).
\item P15 as a psi-field generator is as dead as P1 for the same
  energy-density reason.
\end{itemize}

\subsection{Impact on P3 (precision timing / cavities)}
\label{ssec:P3_impact}

P3 uses precision cavity measurements to detect psi coupling. The
two-component decomposition is neutral for P3: the measurement is in
the $\dpsiEM$ sector, and the two-component split does not change the
$\delta Q/Q$ or $\delta\omega/\omega$ observables.

\subsection{The MOND recovery and laboratory implications}
\label{ssec:MOND_lab}

The recovered MOND formula $a_0 = cH_0\sqrt{\Omega_\Lambda}$ operates
in the $\psigrav$ sector and is irrelevant to laboratory experiments.
However, it restores the astrophysical anchor for $\muG = 0.657$:
MOND galaxy fits constrain $\psigrav$, which through the DFD metric
constrains $\muG$. Since $\Meff = \muG\Vcav/(4\pi G)$, the
astrophysical constraint on $\muG$ propagates to the laboratory
parameters.

\begin{finding}[$\muG$ Remains Well-Constrained]
\label{find:muG_constrained}
The EM-only paradigm shift and two-component decomposition do not
change $\muG = 0.657 \pm 0.006$ (from SPARC galaxy fits). This parameter
enters all P2 sensitivity calculations through $\Meff$. The MOND
recovery via $a_0 = cH_0\sqrt{\Omega_\Lambda}$ ensures that the
galactic rotation curve constraints remain valid, because they depend
on $\psigrav$ (which retains its Mach-integral value
$\psigrav^{\rm cosmo} \approx 0.047$).
\end{finding}

%=============================================================================
\section{W3i6 P2 Summary and Corrections}
\label{sec:P2_summary}
%=============================================================================

\begin{tcolorbox}[colback=blue!3!white, colframe=blue!50!black,
  title=\textbf{W3i6 P2: Summary of Findings and Corrections}]

\begin{enumerate}[label=\textbf{\arabic*.}, itemsep=4pt]

\item \textbf{Passive detection is the only viable path.}
  Active P2 detection (SQUID, MEMS+P11) requires $\Qpsi > 10^{10}$--$10^{13}$,
  which is unmeasured. Passive bounds bypass $\Meff$ and $\Qpsi$ entirely.
  [\textbf{VERIFIED} --- strengthened from W3i5]

\item \textbf{Definitive detection ranking established.}
  LIGO 6th channel ($10^{-19}$) $>$ ESS Lund ($10^{-12}$--$10^{-14}$) $>$
  SQMS next-gen ($10^{-10}$) $>$ MEMS+P11 ($10^{-5}$--$10^{-11}$) $>$
  P1 timing (dead). [\textbf{NEW}]

\item \textbf{Five-step optimal programme defined.}
  (1) LIGO archival, (2) SQMS multi-frequency, (3) ESS Lund,
  (4) $\Qpsi$ indirect measurement, (5) MEMS+P11 conditional prototype.
  [\textbf{NEW}]

\item \textbf{$\Qpsi$ threshold refined.}
  MEMS+P11 requires $\Qpsi \geq 4\ee{10}$ (with 1~$\mu$m NEMS) to
  $4\ee{13}$ (with 10~$\mu$m MEMS). The thermalization upper bound is
  $\sim 10^{45}$, so these values are not excluded. [\textbf{NEW --- CORRECTED}
  from W3i5 estimate]

\item \textbf{GHz MEMS+SQUID coupling is a technology gap.}
  SQUID readout at GHz is not demonstrated. cQED microwave readout is
  the likely alternative, but the psi-mode displacement at the current
  bound is $10^{12}\times$ below $x_{\rm zpf}$. [\textbf{NEW}]

\item \textbf{Two-component decomposition: neutral for P2.}
  The $\psigrav/\dpsiEM$ split does not change P2 sensitivity
  calculations. $\muG$ remains well-constrained by the recovered MOND.
  [\textbf{VERIFIED}]

\item \textbf{Volume Cancellation and Passive Superiority Theorems confirmed.}
  Inherited from W3i5-P2 without modification. These are structural
  results of the DFD framework. [\textbf{VERIFIED}]

\item \textbf{Indirect $\Qpsi$ measurement programme proposed.}
  Three-pronged approach: multi-frequency $Q_0$, temperature-dependent
  $Q_0$, dual-cavity correlation. All executable at SQMS with existing
  or near-term infrastructure. [\textbf{NEW}]

\end{enumerate}
\end{tcolorbox}

%=============================================================================
\section{W3i6 Corrections to Previous Iterations}
\label{sec:corrections}
%=============================================================================

\begin{table}[H]
\centering
\caption{W3i6 corrections and new findings relative to W3i5.}
\label{tab:corrections}
\renewcommand{\arraystretch}{1.2}
\begin{tabular}{@{}p{4.5cm}p{3cm}p{3cm}l@{}}
\toprule
\textbf{Item} & \textbf{W3i5 Value} & \textbf{W3i6 Value} & \textbf{Status} \\
\midrule
T5 impact on P1 & Devastating & Resolved (two-comp.) &
  \textcolor{strongcolor}{\textbf{IMPROVED}} \\
P1 experimental status & Dead & Dead (confirmed) &
  \textcolor{falsified}{\textbf{CONFIRMED DEAD}} \\
P2 strategic recommendation & Deprioritise & Five-step programme &
  \textcolor{strongcolor}{\textbf{NEW}} \\
$\Qpsi$ threshold & $10^{15}$ stated & $4\ee{10}$--$4\ee{13}$ derived &
  \textcolor{strongcolor}{\textbf{CORRECTED}} \\
MEMS+SQUID feasibility & Speculative & GHz gap identified &
  \textcolor{conjcolor}{\textbf{NEW OBSTACLE}} \\
$\muG$ status & Uncertain & Constrained (MOND recovered) &
  \textcolor{strongcolor}{\textbf{IMPROVED}} \\
Definitive signal ranking & Not compiled & Table~\ref{tab:definitive_ranking} &
  \textcolor{strongcolor}{\textbf{NEW}} \\
Indirect $\Qpsi$ measurement & Not proposed & Three-pronged programme &
  \textcolor{strongcolor}{\textbf{NEW}} \\
Two-comp.\ impact on P1 & Not assessed & T5 resolved, no new signals &
  \textcolor{strongcolor}{\textbf{NEW}} \\
\bottomrule
\end{tabular}
\end{table}

%=============================================================================
\section{Open Questions for W3i7}
\label{sec:open_questions}
%=============================================================================

\begin{enumerate}[label=\textbf{OQ-\arabic*.}, itemsep=4pt]

\item \textbf{LIGO 6th channel feasibility}: What matched-filter templates
  are needed for a psi-field signal in LIGO strain data? What EM sources
  provide the strongest psi-correlated modulation? (Requires P6 + P10 input.)

\item \textbf{ESS LLRF noise floor}: What is the realistic systematic
  floor for beam-correlated $Q_0$ modulation at ESS? Can digital LLRF
  provide $\delta Q/Q < 10^{-6}$?

\item \textbf{$\Qpsi$ theoretical prediction}: Can the DFD field equations
  predict $\Qpsi$ from first principles? What dissipation mechanisms
  exist for the psi mode (mode-mixing, nonlinear self-coupling, coupling
  to gravitational radiation)?

\item \textbf{cQED psi-mode readout}: Can circuit QED techniques be
  adapted to read out a psi-mode displacement in a NEMS resonator?
  What is the quantum-limited sensitivity for this transduction chain?

\item \textbf{Two-component decomposition}: formal derivation.
  The $\psi = \psigrav + \dpsiEM$ split was proposed in W3i5 as a
  physical interpretation. Can it be derived from the DFD action
  principle? Is the split unique? Are there cross-terms in the
  action that couple $\psigrav$ and $\dpsiEM$?

\item \textbf{P7 corrected storage time}: With the P7 storage time
  now 1~s (not 18.5~hr), how does this affect the P7--P2 interaction?
  (Deferred to P7 agent.)

\end{enumerate}

%=============================================================================
\section*{Literature References}
%=============================================================================

\begin{thebibliography}{99}

\bibitem{W3i5P1}
DFD Research Programme, Agent~1,
``W3i5 P1: M$_\psi$ Correction Propagation, T1/T5 Impact, Corrected
Smoking-Gun,''
\textit{DFD Research Output}, W3i5\_P1\_update.tex, March 2026.

\bibitem{W3i5P2}
DFD Research Programme, Agent~2,
``W3i5 P2: Complete Reassessment with Corrected $M_{\rm eff}$, Honest
Detection Gap Analysis,''
\textit{DFD Research Output}, W3i5\_P2\_update.tex, March 2026.

\bibitem{W3i5P14P15}
DFD Research Programme, Agent~14/15,
``W3i5 P14+P15: T1 Resolution via EM-Only Sourcing, $M_{\rm eff}$-Corrected
Detection Roadmap,''
\textit{DFD Research Output}, W3i5\_P14\_P15\_update.tex, March 2026.

\bibitem{W3i5CosmoMatSci}
DFD Research Programme, Cosmo+MatSci Agent,
``W3i5: EM-Only Sourcing Resolution of T1--T5, Two-Component $\psi$-Field,''
\textit{DFD Research Output}, W3i5\_Cosmo\_MatSci\_update.tex, March 2026.

\bibitem{W3i5P10P11}
DFD Research Programme, Agent~10/11,
``W3i5 P10+P11: Corrected SQUID-Active, Passive Bounds, LIGO 6th Channel,''
\textit{DFD Research Output}, W3i5\_P10\_P11\_update.tex, March 2026.

\bibitem{SQMS2024}
A.~Romanenko \textit{et al.},
``Ultra-high quality factors in SRF cavities at SQMS,''
\textit{Phys. Rev. Accel. Beams}, 2024.

\bibitem{NEMS-GHz}
M.~Li \textit{et al.},
``GHz-frequency nanoelectromechanical resonators,''
\textit{Nature Nanotechnology}, 2024.

\bibitem{cQED}
A.~Reed \textit{et al.},
``Faithful conversion of propagating quantum information to mechanical motion,''
\textit{Nature Physics}, 2017.

\end{thebibliography}

\end{document}
